% !TeX root = ../教材题库_全部条目_转写工作稿.tex
% Source: XJ-B1-C03-EX-P074-Y0394
% 第3章 / 3.1 函数 / 教材 p.74 / PDF 第81页
% OCR未校对：公式与文字均可能有误；未完成数学TeX转写；禁止直接发给学生
% 原题及相邻原解截图（核对用）：../题目原图/XJ-B1-C03-EX-P074-Y0394-01.webp
% 本文件只在转写完成后作为题干使用；OCR 条目尚不是可靠题干。
\begin{bankitem}{3.1 例题 例题6}{第3章 / 3.1 函数 / 教材 p.74 / PDF 第81页}{XJ-B1-C03-EX-P074-Y0394}{OCR未校对：公式与文字均可能有误；未完成数学TeX转写；禁止直接发给学生}
\begin{ocrdraft}
[教材p.74 / PDF第81页]
\par
0河为推进生态优先、节约集约、绿色低碳发展，推动形成绿色低碳的生产
\par
方式和生活方式，某地采用分段计费的方法计算用户的电费：每月用电量不超过
\par
100kW．h,按0．57元/（kW．h）计费；每月用电量超过100kW．h，其中100kW．h
\par
仍按原标准收费，超过部分按1.5元/()W\ensuremath{\cdot}h）计费．
\par
（1)设月用电工kW\ensuremath{\cdot}h，应交电费3，元，写出3，关于的函数解析式．
\par
（2）小赵家第一季度缴纳的电费情况如下表：
\par
月份
\par
1
\par
2
\par
3
\par
合计
\par
计费金额／元
\par
114
\par
45．6
\par
234．6
\par
问：小赵家第一季度共用电多少？
\par
解
\par
（1）当0主<100时，月电费一月用电量\ensuremath{\times}标准电价，可得
\par
3，一0．57T，
\par
当>100时，月电费一100kW．h的电费十超过100kW．h部分的电费，
\par
可得
\par
3，一0．57\ensuremath{\times}100十1.5\ensuremath{\times}（工一100）=1.5一93．
\par
0，57工，
\par
xeCO，10的，
\par
所以
\par
1.5x一93，xe（100，十亠）．
\par
（2）由（1）可知，当电费不超过57元时，说明月用电量不超过100kW\ensuremath{\cdot}h；当
\par
电费超过57元时，说明月用电量超过100kW．h.因此，用电量应使用函数的不同
\par
关系式来计算．
\par
因为1月份、2月份电费超过57元，所以按第二个函数关系式计算，即1.5x—
\par
93一114，5x一93一75，分别算出1月份用电138kW．h，2月份用电112kW．h；
\par
而3月份电费不超过57元，按第一个函数关系式计算，有0，5五=45．6，算出3月
\par
份用电80kW．h．
\par
因此，小赵家第一季度共用电330kW．h．
\par
一般地，如果自变量在定义域的不同取值范围内时，函数由不同的解析式给出，
\par
这种函数叫作分段函数
\par
例6中的函数就是一个分段函数．
\end{ocrdraft}
\end{bankitem}
